\chapter{HƯỚNG PHÁT TRIỂN}
\label{ch:huong-phat-trien}

Các hướng phát triển được đề xuất sau khi hoàn thiện mức sàn gồm đối chứng cơ sở, loại bỏ đặc trưng,
cuốn chiếu theo thời gian, kiểm tra mốc cắt và tiền xử lý chỉ học từ phần huấn luyện. Thứ tự này có
ý nghĩa phương pháp: chỉ nên tăng độ phức tạp sau khi biết quy trình hiện tại hợp lệ và
kết quả có thể tái lập. Nếu chưa khóa được dữ liệu và giao thức, việc thêm mô hình mới
khó cho biết cải thiện đến từ kiến trúc, từ dữ liệu mới hay từ một thay đổi không được
kiểm soát.

Các hướng dưới đây không phải là kết quả đã được thực hiện trong mức sàn. Mỗi hướng
được xem như một câu hỏi mở, cần có dữ liệu, cấu hình, phép đối chứng và tiêu chí đánh
giá riêng. Thứ tự ưu tiên bắt đầu từ chất lượng dữ liệu và độ tin cậy của giao thức,
sau đó mới mở rộng biểu diễn, mô hình, phạm vi thị trường và cuối cùng là giá trị kinh
tế. Cách sắp xếp này giúp báo cáo không nhầm kế hoạch phát triển với bằng chứng thực
nghiệm đã đo.

\section{Hạng mục của đề cương được chuyển sang hướng phát triển}
\label{sec:hpt-chuyen}

Ba hạng mục nêu trong đề cương không được thực hiện trong phạm vi thời gian của đồ án. Việc
chuyển chúng sang chương này là một quyết định về thứ tự ưu tiên, được ghi lại tường minh
thay vì để trống:

\begin{enumerate}[nosep]
  \item \textbf{Temporal Fusion Transformer.} Chương~\ref{ch:phuong-phap} đã đặt TFT ở vị
        trí mở rộng có điều kiện, chỉ thực hiện sau khi mức sàn đạt cổng chất lượng. Kết quả
        ở Chương~\ref{ch:ket-qua} cho thấy đóng góp thông tin của cảm xúc không khác 0, nên
        thay LSTM bằng một kiến trúc phức tạp hơn khó làm thay đổi kết luận về câu hỏi
        trung tâm. Ưu tiên đúng là nâng chất lượng tín hiệu đầu vào trước, rồi mới tăng độ
        phức tạp của bộ dự báo. Khi thực hiện, TFT cần được cấu hình cho phân loại ba lớp và
        chạy trên đúng tập giữ lại; phần quan trọng của biến và cơ chế chú ý chỉ được diễn
        giải như mô tả phụ thuộc mô hình, không phải bằng chứng nhân quả.
  \item \textbf{Biến thể PhoBERT kết hợp CNN hoặc BiLSTM.} Xác thực chéo mới trên 1.306
        mẫu đã cho mức đo lớp NEGATIVE tốt hơn, nhưng chưa truyền bộ sinh mới sang nhánh
        dự báo. Trước hết cần tinh chỉnh lại hoặc tổ hợp các điểm kiểm hiện có, chấm lại
        kho tin và chạy lại đối chứng; sau đó mới so sánh các biến thể kiến trúc ở
        Mục~\ref{sec:kq-mornhan}.
  \item \textbf{Hệ thống demo trực quan hóa.} Toàn bộ quy trình hiện được điều khiển qua
        giao diện dòng lệnh và ghi artifact ra tệp, đủ cho mục đích nghiên cứu và tái lập.
        Một giao diện trực quan hóa sẽ hữu ích cho việc thăm dò dữ liệu, nhưng không bổ sung
        bằng chứng cho câu hỏi nghiên cứu, nên được xếp sau các hạng mục làm tăng độ tin cậy
        của kết quả.
\end{enumerate}

Ưu tiên cao nhất trong số các hướng của chương này là truyền bộ sinh cảm xúc đã đánh giá
qua toàn bộ chuỗi dự báo: tinh chỉnh lại hoặc tổ hợp các điểm kiểm, chấm lại kho tin theo
mốc cắt thông tin, rồi chạy lại dự báo và đối chứng. Chỉ chuỗi thực nghiệm hoàn chỉnh này
mới có thể tách tín hiệu cảm xúc yếu khỏi ảnh hưởng của bộ sinh lịch sử.

\section{Củng cố dữ liệu và giao thức point-in-time}

Ưu tiên đầu tiên là nâng chất lượng dữ liệu. Có thể mở rộng nguồn tin ngoài Vietstock,
chẳng hạn CafeF, VnEconomy hoặc Stockbiz, nhưng mỗi nguồn mới cần được kiểm tra về độ
phủ theo mã, mốc thời gian, nội dung, trùng lặp, điều kiện truy cập và khả năng tái lập. Cần
so sánh không chỉ số lượng bài viết mà cả tỷ lệ bài có thể ánh xạ được với mã và phiên
giao dịch. Một thông cáo được đăng lại ở nhiều nguồn phải được nhận diện để không làm
tăng giả tạo trọng số của cùng một sự kiện. Bản ghi thiếu thời gian hoặc không xác định
được mã nên được loại khỏi tập chính và ghi vào thống kê độ bao phủ, thay vì lấp giá trị
thiếu bằng một thời điểm được suy đoán.

Mốc cắt 15 giờ có thể được kiểm tra độ nhạy bằng một thực nghiệm riêng. Có thể so
sánh tin trước giờ đóng cửa với tin sau giờ đóng cửa được chuyển sang phiên kế tiếp,
nhưng mỗi mốc cắt phải tạo ra bảng đặc trưng và tập giữ lại riêng. Không được chọn mốc cắt
theo điểm kiểm thử tốt nhất mà không ghi số lần thử. Tương tự, có thể so sánh ngưỡng phân vị
với ngưỡng cố định, nhưng cách xác định ngưỡng phải tuân thủ tập huấn luyện-only ở mọi trường
hợp. Việc kiểm tra độ nhạy chỉ có ý nghĩa nếu các mô hình trong các cấu hình vẫn dùng
cùng ngày đánh giá và cùng cách tính chỉ số.

Một hướng khác là làm giàu annotation. Bộ nhãn có thể bổ sung cường độ, chủ thể, loại
sự kiện và mức độ chắc chắn nếu có hướng dẫn cùng đánh giá độ đồng thuận. Nhãn sự kiện
chi tiết có thể bổ sung cho cảm xúc vì một bài viết đôi khi truyền đạt loại sự kiện
quan trọng mà ba lớp phân cực không thể hiện \cite{Chen2019_191005078v1}. Khi mở rộng,
cần phân biệt nhãn tác động đến doanh nghiệp với nhãn tác động đến toàn ngành hoặc chỉ
số. Cũng cần giữ mẫu bất đồng để phân tích ranh giới khó, không loại bỏ toàn bộ những
mẫu này chỉ vì chúng làm điểm đánh giá giảm. Đây là mở rộng câu hỏi nghiên cứu, không
đưa vào so sánh chính nếu làm thay đổi mục tiêu ban đầu.

\section{Cơ chế hợp nhất nâng cao}

Mức sàn dùng phép nối vì dễ kiểm tra. Khi dữ liệu đủ lớn, có thể thử có cổng hợp nhất, trong
đó một cổng học cách điều chỉnh mức đóng góp của nhánh cảm xúc theo trạng thái giá,
độ biến động và số tin. Một hướng khác là chú ý chéo hai chiều, cho phép các bước
giá truy vấn chuỗi cảm xúc và ngược lại. Các mô hình Transformer đa phương thức cho
thấy hợp nhất động có thể giữ được tương tác phong phú hơn phép nối trực tiếp
\cite{Omranpour2024_241210540v1,Elahi2024_241101368v1}.

Hợp nhất động có nhiều tham số hơn và có thể học nhiễu khi số mã hoặc số ngày có tin còn
ít. Mỗi biến thể phải được so sánh bằng cùng tập giữ lại, ngân sách thử, loại bỏ đặc trưng và chỉ số.
Cần báo cáo số tham số, thời gian huấn luyện, độ nhạy theo hạt giống, số ngày có tin thực sự
được sử dụng và khoảng tin cậy. Nếu một biến thể chỉ tốt ở một phần chia hoặc chỉ tốt ở nhóm
có rất nhiều tin, không nên gọi đó là cải tiến ổn định trước khi kiểm tra kích thước mẫu
và các yếu tố đồng thời.

TFT cũng có thể được mở rộng thành mô hình đa nhiệm, trong đó một tầng phân loại dự báo lớp xu
hướng và một tầng phân loại dự báo độ biến động hoặc xác suất không có tin. Hướng này chỉ phù hợp
khi nhãn phụ có cơ sở, hàm mất mát được cân bằng và không làm rò rỉ mục tiêu. Cần có thí nghiệm
đối chứng với TFT chỉ có mục tiêu chính để biết tầng phân loại phụ có thực sự giúp học tốt hơn hay
chỉ làm tăng số tham số. tầm quan trọng của biến và chú ý vẫn phải được diễn giải như
tín hiệu nội tại của mô hình, không phải quan hệ nhân quả.

\section{Mở rộng phạm vi thị trường và trôi dạt khái niệm}

Sau khi quy trình ổn định trên rổ mã ban đầu, có thể mở rộng sang toàn bộ VN30, HNX,
UPCOM hoặc các nhóm ngành khác. Việc mở rộng nên thực hiện từng bước. Trước hết kiểm
tra độ bao phủ, tính nhất quán của lịch giao dịch và sự khác nhau về nguồn tin, sau đó
đánh giá khả năng tổng quát theo mã. Không nên gộp mọi mã vào một điểm số duy nhất nếu
mật độ tin, thanh khoản hoặc số phiên hợp lệ giữa các mã chênh lệch lớn. Mỗi lần thêm
mã cần ghi ngày bắt đầu, lý do đưa vào, số mẫu bị loại và ảnh hưởng của thay đổi đó đến
phân bố nhãn.

Các hướng quốc tế như HIST và REST khai thác quan hệ giữa cổ phiếu, còn DoubleAdapt
nhắm đến trôi dạt khái niệm trong dữ liệu liên tục \cite{Xu2021_211013716v2,
Xu2021_210207372v2,Zhao2023_230609862v3}. Đồ thị ngành hoặc đồ thị suy ra từ tương
quan có thể được thử nghiệm, nhưng phải làm rõ nguồn quan hệ, thời điểm xây dựng và
nguy cơ dùng thông tin tương lai. Nếu đồ thị được học từ dữ liệu, nó phải được dựng
trong tập huấn luyện của từng phần chia, không dựng một lần từ toàn bộ giai đoạn rồi dùng cho kiểm thử.
MTMD, DishFT-GNN và S$^{3}$G có thể cung cấp gợi ý về bộ nhớ đa quy mô, chưng cất mô hình
và đồ thị không gian trạng thái, nhưng không phải mục tiêu bắt buộc
\cite{Wang2022_221208656v2,Liu2025_250210776v1,Lu2026_260324236v1}.

Một hướng thực tế hơn là kiểm tra trôi dạt khái niệm bằng cách huấn luyện trên các cửa sổ
khác nhau, theo dõi chênh lệch giữa tập huấn luyện và kiểm thử, rồi cập nhật điểm kiểm theo lịch đã
khóa. Có thể so sánh cửa sổ mở rộng với cửa sổ trượt để xem lịch sử xa còn hữu ích hay
không. Nếu hệ thống cập nhật trực tuyến, cần ghi phiên bản dữ liệu, thời điểm cập nhật,
điểm kiểm được dùng và khả năng rollback. Không nên cập nhật mô hình sau mỗi phiên mà
không có đánh giá riêng, vì cách làm đó khó phân biệt thích nghi với quá khớp.

\section{Đánh giá giá trị kinh tế}

Dự báo đúng lớp chưa đủ để kết luận tín hiệu có giá trị giao dịch. Hướng phát triển
quan trọng là xây dựng kiểm thử quá khứ với quy tắc rõ về thời điểm vào lệnh, giá khớp, phí,
trượt giá, tỷ trọng, thanh khoản và giới hạn vị thế. Quy tắc chuyển xác suất ba lớp
thành vị thế cũng phải được khóa trước khi xem kết quả. Các chỉ số có thể gồm lợi suất tích lũy,
lợi suất năm hóa, tỷ số Sharpe, mức suy giảm tối đa, vòng quay danh mục và độ ổn định theo
giai đoạn.

Kiểm thử quá khứ phải dùng dự đoán ngoài mẫu từ cuốn chiếu theo thời gian. Không dùng dự đoán của tập giữ lại để
điều chỉnh chiến lược rồi báo cáo lại trên chính tập giữ lại. Nếu so sánh nhiều chiến lược,
số lần thử phải được ghi lại và kết luận cần thận trọng trước kiểm định nhiều lần. Các
nghiên cứu dùng mô phỏng giao dịch cho thấy đánh giá kinh tế có thể bổ sung cho chỉ số
phân loại, nhưng kết quả vẫn phụ thuộc giả định mô phỏng và thị trường
\cite{Chen2021_210708721v1,Feng2018_181009936v2}.

Kiểm thử quá khứ không thuộc mức sàn vì nó mở rộng câu hỏi từ dự báo sang quyết định giao dịch.
Khi thực hiện, chương trình sinh dự đoán và chương trình mô phỏng nên tách nhau để thay
đổi chiến lược không làm thay đổi dữ liệu dự báo. Kết quả kinh tế cũng cần được phân
tách theo giai đoạn, mã và mức chi phí giả định; một lợi nhuận tích lũy duy nhất không
đủ để mô tả độ bền của tín hiệu.

\section{Giải thích và kiểm tra độ bền}

Có thể mở rộng giải thích bằng loại bỏ đặc trưng, che khuất, tầm quan trọng hoán vị,
Integrated Gradients cho nhánh văn bản hoặc phân tích token quan trọng. Mỗi phương pháp
có giả định riêng và có thể cho kết quả khác nhau. Vì vậy, chúng nên được xem là bằng
chứng bổ sung, được đối chiếu với kiểm tra loại biến, kết quả theo nhóm và trường hợp
cụ thể. Phân tích token được chú ý không đồng nghĩa với việc token đó gây ra biến động
giá; kết luận vẫn phải giới hạn ở cách mô hình sử dụng đầu vào.

Phân tích độ bền nên bao gồm thay đổi hạt giống, độ dài cửa sổ, mốc cắt, ngưỡng nhãn và tập
mã. Có thể phân tầng theo ngày có sự kiện, số tin, độ biến động và hướng thị trường.
Nhóm có quá ít quan sát cần được đánh dấu là không đủ để kết luận. Nếu chênh lệch
cảm xúc chỉ xuất hiện khi có nhiều tin, cần kiểm tra xem đó là tín hiệu nội dung hay
chỉ là ảnh hưởng của cờ \texttt{has\_news}, độ phủ nguồn và kích thước mẫu. Các phân
tầng phải được xác định trước hoặc ghi rõ số lần thử để tránh chỉ báo cáo nhóm có kết
quả thuận lợi.

Một hướng khác là hiệu chỉnh xác suất. Bên cạnh lớp argmax, có thể đo Brier score,
Expected Calibration Error hoặc hiệu chỉnh trên xác thực. Điều này hữu ích nếu hệ
thống được dùng để xếp hạng độ tin cậy, nhưng không thay thế F1 vĩ mô trong giao thức
chính.

\section{Mở rộng mô hình ngôn ngữ}

Khi kho ngữ liệu đủ lớn, có thể so sánh PhoBERT-base với mô hình nền tiếng Việt khác hoặc tiếp
tục tiền huấn luyện trên văn bản tài chính không nhãn. Mọi bước tiền huấn luyện phải lưu phiên bản
kho ngữ liệu và kiểm tra không đưa văn bản có nhãn hoặc thông tin tương lai vào đánh giá.
PhoBERT phù hợp với tiếng Việt, nhưng điểm kiểm lớn hơn không mặc nhiên tốt hơn
\cite{Nguyen2020_200300744v3}. Thí nghiệm này cần tách ảnh hưởng của tiền huấn luyện khỏi
ảnh hưởng của tầng phân loại, bộ tách từ, độ dài đầu vào và cách chia dữ liệu.

Có thể thử biểu diễn kết hợp tiêu đề và phần mở đầu, trong đó một bộ mã hóa xử lý tiêu
đề và một bộ mã hóa xử lý nội dung. Hướng này cần đánh giá riêng ảnh hưởng của độ dài,
truncation, quảng cáo và bản tin trùng lặp. FinGPT hoặc hệ thống hợp nhất qua LLM có
thể làm đối chứng cơ sở tham khảo, nhưng phải ghi rõ phiên bản, chi phí, câu lệnh hoặc cấu hình
và khả năng tái lập; chúng không nên thay PhoBERT trong quy trình sàn
\cite{Liang2024_241210823v2,Siala2026_260200086v3}.

\section{Đóng gói và tái lập hệ thống}

Một phiên bản phát triển hoàn chỉnh nên có các mô-đun độc lập cho thu thập, chuẩn hóa
và loại trùng, ánh xạ theo phiên, suy luận cảm xúc, tạo bảng đặc trưng, huấn luyện,
đánh giá và trực quan hóa. Các mô-đun cần chia sẻ lược đồ dữ liệu và bản kê, không đọc
các file tạm không có phiên bản. Mỗi đầu ra trung gian nên có thông tin về nguồn, thời
điểm tạo, cấu hình và hash để có thể truy lại bước đã sinh ra nó.

Hệ thống demo có thể hiển thị tin đã thu thập, xác suất ba lớp, số tin theo mã, đặc
trưng kỹ thuật, lớp dự báo và cảnh báo dữ liệu thiếu. Giao diện chỉ nên trình bày thời
điểm tạo dự đoán, điểm kiểm và độ tin cậy; không gọi kết quả là khuyến nghị mua bán.
Mỗi bản ghi demo phải cho biết dữ liệu được cập nhật đến thời điểm nào, mốc cắt nào được
áp dụng và mô hình nào tạo ra kết quả để tránh nhầm giữa hồi cứu và dự báo thực.

Để tăng khả năng tái lập, cần công bố file cấu hình, hạt giống, mã băm dữ liệu, phiên bản
điểm kiểm, nhật ký chia tập, số lần thử và cách dựng môi trường. Nếu dữ liệu gốc không
thể phân phối do điều khoản nguồn, có thể công bố lược đồ, hash, quy trình và bản kê đã
ẩn nội dung nhạy cảm, kèm giới hạn truy cập. Khi nguồn tin thay đổi cấu trúc hoặc xóa
bài, cần lưu trạng thái lỗi và thời điểm thu thập để người khác hiểu vì sao việc chạy
lại có thể không tạo cùng số bản ghi.

\section{Lộ trình ưu tiên}

Lộ trình phát triển được đề xuất như sau:

\begin{enumerate}[nosep]
  \item Hoàn thiện cổng chất lượng dữ liệu, mốc cắt, nhãn và kiểm tra rò rỉ;
  \item Khóa đối chứng cơ sở, LSTM hai nhánh, loại bỏ đặc trưng và khoảng tin cậy theo lấy mẫu lặp trên tập giữ lại;
  \item Kiểm tra độ bền theo hạt giống, cửa sổ, mã và độ bao phủ tin;
  \item Thử TFT phân loại nếu còn dữ liệu, thời gian và tài nguyên;
  \item Bổ sung hợp nhất động hoặc tín hiệu ở mức sự kiện khi giao thức đã ổn định;
  \item Mở rộng mã, nguồn và trôi dạt khái niệm sau khi kiểm tra khả năng tái lập;
  \item Cuối cùng mới triển khai kiểm thử quá khứ và các chỉ số giá trị kinh tế.
\end{enumerate}

Thứ tự này bảo vệ mục tiêu cốt lõi của đồ án. Các hướng mở rộng chỉ có ý nghĩa khi
được đánh giá trên nền dữ liệu sạch, có đối chứng và có thể truy nguyên.
